\documentclass[border=14pt,varwidth=18.5cm]{standalone}
\usepackage{amsmath,amssymb}
\usepackage{mathtools}
\usepackage{empheq}
\usepackage[T1]{fontenc}
\usepackage{lmodern}
\usepackage{xcolor}
\usepackage{booktabs}
\usepackage{array}
\usepackage[american,siunitx]{circuitikz}
\definecolor{head}{HTML}{0072B2}
\definecolor{hl}{HTML}{FFF59D}
\definecolor{note}{HTML}{555555}
\newcommand{\sheettitle}[1]{{\Large\bfseries\color{head}#1}\par\smallskip
  {\small\color{note}Jeremia Christ Immanuel Manalu - 5023231017 - Biocybernetics (EB234701), Class Q}\par\medskip}
\newcommand{\sect}[1]{\par\medskip{\large\bfseries\color{head}#1}\par\smallskip}
\newcommand{\hlbox}[1]{\colorbox{hl}{$\displaystyle #1$}}
\newcommand{\nomtable}[1]{\par\medskip{\large\bfseries\color{head}Notation (in order of appearance)}\par\smallskip
  \renewcommand{\arraystretch}{1.15}%
  \begin{tabular}{>{$}l<{$} p{10.6cm} l}
  \toprule
  \textrm{\textbf{Symbol}} & \textbf{Meaning} & \textbf{Unit} \\
  \midrule
  #1
  \bottomrule
  \end{tabular}}
\setlength{\parindent}{0pt}

\renewcommand{\theequation}{6.\arabic{equation}}
\begin{document}
\sheettitle{Analysis Definitions: Spike, Inter-Pulse Interval, Refractory Threshold and Conduction Velocity}

\sect{Action potential (spike) detection}
\begin{align}
t_{c,i} &= \text{instant of the $i$-th upward crossing of } V_m = 0\ \text{mV},\qquad
V_{peak,i} = \max_{t\in[t_{c,i},\,t_{d,i}]} V_m(t),\quad t_{p,i}=\arg\max
\end{align}
where $t_{d,i}$ is the next downward crossing. Graded responses that stay below 0~mV are not counted.

\sect{Inter-pulse interval and firing frequency (tab 4)}
\begin{align}
\mathrm{ISI}_i &= t_{p,i}-t_{p,i-1},\qquad
\hlbox{\overline{\mathrm{ISI}}(J_{in}) = \frac{1}{N_s}\sum_{t_{p,i-1}\ge t_{dis}}\mathrm{ISI}_i,\qquad f = \frac{1000}{\overline{\mathrm{ISI}}}\ \text{Hz}}
\end{align}
A constant $J_{in}$ is applied from $t=0$; spikes before $t_{dis}$ (transient) are discarded. Larger $|J_{in}|$ charges $C_m$ faster
after each spike, so the threshold is reached earlier and $\overline{\mathrm{ISI}}$ decreases: the input level programs the inter-pulse interval.

\sect{Refractory period (paired-pulse protocol, tab 3)}
\begin{align}
J_{th,0} &= \min\{\,|J|:\ \text{one pulse of width } w \text{ evokes an AP}\,\}\qquad\text{(rested threshold)} \\
J_{th}(\Delta T) &= \min\{\,|J_2|:\ \text{the test pulse at } t_{on}+\Delta T \text{ evokes a second AP}\,\} \\
&\hlbox{\text{absolute: } J_{th}(\Delta T)>|J_2|_{max};\qquad \text{relative: } J_{th,0}<J_{th}(\Delta T)\ \text{and}\ J_{th}\le 1.05\,J_{th,0}\ \text{marks its end}}
\end{align}
Both thresholds are found by bisection on $[0,\ |J_2|_{max}]$ with 16 halvings (resolution $|J_2|_{max}/2^{16}$).

\sect{Conduction velocity and passive length constant (tab 5)}
\begin{align}
&\hlbox{\theta = \frac{x_{75\%}-x_{25\%}}{t_c(x_{75\%})-t_c(x_{25\%})}\cdot 10\ \ \text{m/s}\quad(\text{cm/ms}\to\text{m/s})} \\
\lambda &= \sqrt{\frac{1}{2\pi a\,g_m\,(r_o+r_i)}},\qquad \theta \propto \sqrt{a}\ \ \text{(HH 1952, eq. 34)}
\end{align}
The crossing time $t_c(x_j)$ is linearly interpolated between the two time levels that bracket $V_m=0$.

\nomtable{
t_{c,i},t_{d,i} & upward and downward 0 mV crossing times of spike $i$ & ms \\
V_m & absolute membrane potential & mV \\
V_{peak,i},t_{p,i} & peak value and peak time of spike $i$ & mV, ms \\
\mathrm{ISI}_i & inter-spike (inter-pulse) interval & ms \\
\overline{\mathrm{ISI}} & mean steady interval for a given $J_{in}$ & ms \\
N_s & number of intervals used in the mean & - \\
t_{dis} & discarded transient time & ms \\
f & firing frequency & Hz \\
J_{in} & constant stimulus current density & $\mu$A/cm$^2$ \\
J_{th,0} & rested single-pulse threshold & $\mu$A/cm$^2$ \\
w & pulse width & ms \\
J_{th},\ \Delta T & test-pulse threshold and onset-to-onset interval & $\mu$A/cm$^2$, ms \\
J_2,\ t_{on} & test-pulse amplitude and conditioning-pulse onset & $\mu$A/cm$^2$, ms \\
|J_2|_{max} & largest test amplitude searched & $\mu$A/cm$^2$ \\
\theta & conduction velocity & m/s \\
x_{25\%},x_{75\%} & nodes at 25 \% and 75 \% of the axon length & cm \\
\lambda & passive length constant & cm \\
a & axon radius & cm \\
g_m,r_o,r_i & passive conductance and axial resistances (sheet 4) & mS/cm$^2$, $\Omega$/cm \\
}
\end{document}
